%! Quadratic Systems \& Modeling — Unit 2, Lesson 2.5 — Lesson Plan
% SLIDES-FIRST plan (course shape 2026-09-25, engagement rules 2026-09-27, condensed Unit 2
% 2026-09-29): the deck carries the lesson — warm-up, three rounds of definition → example →
% ladder → Show me / Find the mistake, a Final Round, and the homework (DeltaMath online) started
% in class. Students take their notes on the printed slide handout. Teacher-facing; never handed out.
\documentclass[10pt]{article}
\usepackage{algebra2-article}
\usepackage{algebra2-boxes}
\usepackage{graphicx}   % -article does NOT load graphicx; needed for thumbnails/screenshots
\graphicspath{{images/}}
\newcommand{\TallMath}[1]{$\displaystyle #1\rule[-1.4em]{0pt}{3.2em}$}
\newcommand{\CourseName}{Algebra 2}
\newcommand{\MeetingLength}{60 minutes}
\newcommand{\UnitNumberName}{Unit 2: Quadratic Functions \quad}
\newcommand{\LessonNumberName}{Lesson 2.5: Quadratic Systems \& Modeling}
\newcommand{\rung}[1]{{\color{goldacc}\ifcase#1\or$\star$\or$\star\star$\or$\star\star\star$\fi}}

\begin{document}
\begin{center}
    {\Large\bfseries \CourseName \\
    {\small\bfseries \UnitNumberName \LessonNumberName}}
\end{center}

\begin{tcolorbox}[colback=forestbg, colframe=forest, boxrule=0.9pt, arc=2mm,
    left=2mm, right=2mm, top=1.5mm, bottom=1.5mm]
    \textbf{Primary Objective:} Students solve a linear--quadratic system by substitution and
    give each solution as an ordered pair (an intersection point), predict the number of
    solutions from the discriminant of the collapsed equation, and answer questions about a
    quadratic model by choosing the right feature --- the $y$-intercept (start), a positive zero
    (when it reaches $0$), or the vertex, whose \emph{output} is the maximum or minimum value and
    whose \emph{input} is when or where it happens. Every equation solved today is a 2.2--2.4
    equation; the new work is setting it up and reading the answer.
    \\[2pt]
    \textbf{Standards (2023 VA SOL):} \textbf{A2.EI.3a} (create a linear--quadratic system to
    model a situation), \textbf{A2.EI.3b} (determine the number of solutions), \textbf{A2.EI.3c}
    (solve a linear--quadratic system algebraically and graphically; quadratic--quadratic touched
    on a \rung3), \textbf{A2.EI.3d} (verify solutions and interpret them), \textbf{A2.EI.2a}
    (create a quadratic model of a contextual situation), \textbf{A2.EI.2d} (verify and interpret
    solutions in context), \textbf{A2.F.2d} (identify and interpret the maximum/minimum value
    from the vertex).
    \\[2pt]
    \textbf{Lesson model:} \emph{Slides first.} The deck is the lesson and its printed handout is
    the student's notes, and students are \emph{solving the whole hour}: a warm-up on screen as
    they walk in; three rounds of a short \textbf{definition} and \textbf{example} ($\le$3 min), a
    \textbf{ladder} of three problems (\rung1\ \rung2\ \rung3), and a \textbf{Show me} or
    \textbf{Find the mistake}; a self-scored \textbf{Final Round}; and the \textbf{homework}
    started in class. \textbf{Homework source:} DeltaMath online, set
    \textbf{2.5 Systems \& Modeling} --- not in the packet.
\end{tcolorbox}

\renewcommand{\arraystretch}{1.4}
\begin{skillbox}[Priority Ideas \& Skills]{goldbox}
\begin{tabularx}{\linewidth}{@{}X|X@{}}
\textbf{Priority Ideas \& Skills} & \textbf{Key Understandings (the ``why'')} \\[2pt]
\hline
\vspace{-6pt}
\begin{itemize}[leftmargin=1.1em, nosep, after=\vspace{2pt}]
  \item Solve a linear--quadratic system by \textbf{substitution}: set the $y$'s equal, solve,
        back-substitute into the line, check in both.
  \item Count solutions ($2$ / $1$ / $0$) from the \textbf{discriminant} of the collapsed
        equation, and read them off a graph (cross / touch / miss).
  \item In a model, match the question to the feature: $y$-intercept, positive zero, vertex.
  \item Find the vertex by $t=-b/2a$, then evaluate for the max/min value; answer with units.
\end{itemize}
&
\vspace{-6pt}
\begin{itemize}[leftmargin=1.1em, nosep, after=\vspace{2pt}]
  \item A solution is a \emph{point} on both graphs --- so it needs both coordinates.
  \item Substitution collapses two equations into one quadratic, so every 2.2--2.4 tool applies.
  \item The collapsed equation's $D$ counts the intersections before any solving.
  \item \textbf{Target misconception: the input is not the output.} Asked ``how high?'',
        students answer with the vertex's \emph{input} ($1$ second) instead of its
        \emph{output} ($64$ feet) --- or grab the zero or the starting height. The maximum
        \emph{value} is an output; \emph{when} it happens is the input.
\end{itemize}
\end{tabularx}
\end{skillbox}

\begin{skillbox}[{Vocabulary, Concepts, \& Theorems --- one definition frame each}]{sky}
\renewcommand{\arraystretch}{1.3}
\begin{tabularx}{\linewidth}{@{}l X@{}}
\textbf{Solution of a system} & A \textbf{system} is two or more equations taken together. A
    \textbf{solution} is an ordered pair $(x,y)$ that makes \emph{both} equations true --- a point
    where the graphs \textbf{intersect}. \\
\textbf{Solving by substitution} & Both equations say $y=\ldots$: set the right sides equal; move
    everything to one side and solve the quadratic; put each $x$ into the \emph{line} to get its
    $y$; check each point in both equations. \\
\textbf{Number of solutions} & A line and a parabola can \textbf{cross} at $2$ points,
    \textbf{touch} at $1$, or \textbf{miss} ($0$). The collapsed equation's
    \textbf{discriminant} $D=b^2-4ac$ gives $D>0$: $2$; $D=0$: $1$; $D<0$: $0$. \\
\textbf{Feature $\to$ question} & \textbf{$y$-intercept} $h(0)$: the starting value.
    \textbf{Positive zero}: when it lands (throw out $t<0$). \textbf{Vertex} input $-b/2a$: when
    / where it is highest. \textbf{Vertex} output: how high / how much. \\
\textbf{Maximum or minimum value} & The vertex's \textbf{output}. The vertex's \textbf{input}
    says \emph{when} or \emph{where} it happens. \\
\end{tabularx}
\end{skillbox}

\begin{fixedskillbox}[Lesson at a Glance --- \MeetingLength]{forestbg}
\renewcommand{\arraystretch}{1.25}
\begin{tabularx}{\linewidth}{@{}l c X X@{}}
\textbf{Phase} & \textbf{Min} & \textbf{Students} & \textbf{Teacher} \\
\hline
Warm-Up & 5 & Three quick wins, on screen as they walk in; score themselves & Stand at the door; reveal at 3 min; land ``where do the parabola and the line meet?'' \\
Lesson & 40 & Three rounds (13 / 13 / 14): watch (2 min), \emph{climb the ladder} \rung1$\to$\rung2$\to$\rung3, then \emph{Find the mistake} or \emph{Show me} fingers & $\le$3 min talk a round; circulate the whole ladder; reveal when most have \rung2; read the fingers \\
Final Round & 5 & Four mixed problems + bonus, alone; score out of 4 & Timer on; collect the scores on fingers \\
Homework & 10 & Start the DeltaMath set online, alone & Launch problem~1, then circulate \\
\end{tabularx}
\end{fixedskillbox}

\begin{skillbox}[Engagement --- how this hour keeps everyone solving]{goldbox}
\textbf{Nobody waits:} every practice frame is a ladder --- \rung1\ everyone (a quick win),
\rung2\ everyone, \rung3\ early finishers. \textbf{Everyone commits, safely:} \emph{Show me} is
1--4 fingers on the count of three, all at once. \textbf{Critique a stranger:} Round~1's
\emph{Find the mistake} puts the missing $y$-values on an anonymous Student A. \textbf{Their
world:} the volleyball serve and the campus-garden fence on Round~3's ladder, and the art club's
hoodie sale as the bonus. \textbf{Short clock:} every practice frame carries its minutes. Praise
the climb more than the answer.
\end{skillbox}

\begin{skillbox}[Warm-Up (5 min, on the slides as they walk in)]{forestbg}
Three quick wins; the answers are one click. \textbf{(1)} $x=2$ in $y=x^2-2x-3$: $-3$ ---
substitution. \textbf{(2)} the vertex of $y=x^2-4x+1$: $x=-b/2a=2$, $(2,-3)$ --- the engine of
Round~3. \textbf{(3)} $x^2-3x=0$: $x=0$ or $3$ --- exactly the equation Example~1 collapses into.
The reveal ends on \emph{Hold on to this}: where do $y=x^2-2x-3$ and $y=x-3$ meet? That is
Round~1. Then the targets frame (30 seconds).
\end{skillbox}

\begin{skillbox}[The Lesson --- three rounds of definition, example, ladder, check (40 min)]{forestbg}
\begin{multicols}{2}
\textbf{Round 1 --- linear--quadratic systems (13 min).} Definition with the pre-drawn graph (the
two red points) and Example~1 in 3 minutes: $x^2-2x-3=x-3\Rightarrow x^2-3x=0$, so $(0,-3)$ and
$(3,0)$ --- point at the warm-up's item 3. Ladder (5 min): \rung1\ $(2,4)$, $(-1,1)$;
\rung2\ $(1,6)$, $(-3,-2)$; \rung3\ $y=3-x$ first, $(2,1)$, $(-3,6)$ --- watch for $y$ taken
from the wrong $x$. Find the mistake: Student A stopped at $x=3$, $x=-1$; the fix is
$(3,10)$, $(-1,-6)$. Ask ``what would you tell Student A?''

\textbf{Round 2 --- how many solutions (13 min).} Definition (cross / touch / miss on one graph)
and Example~2's table ($D=36$, $0$, $-8$) in 3 minutes; land ``$D$ of the \emph{collapsed}
equation.'' Ladder (4 min): \rung1\ $D=16$: $2$; \rung2\ $D=0$: $1$, at $(1,2)$;
\rung3\ two parabolas, $2x^2-8=0$, $D=64$: $2$. Show me ($y=x^2-4x+7$, $y=2x-2$): answer
\textbf{2}, one ($x^2-6x+9=0$, $D=0$). \textbf{1} = took $D$ of the parabola alone ($-12$);
\textbf{3} = ``a quadratic always has two''; \textbf{4} = does not trust $D$.

\columnbreak
\textbf{Round 3 --- modeling: the crux (14 min).} Definition (the feature-to-question table
beside the graph) and Example~3 ($h=-16t^2+64t+80$: starts $80$ ft, highest at $2$\,s, $144$
ft, lands $5$\,s, reject $-1$) in 3 minutes --- say ``input'' and ``output'' aloud every time.
Ladder (5 min): \rung1\ $6$ ft; \rung2\ $0.5$\,s, $10$ ft; \rung3\ the garden,
$A=-2x^2+40x$: width $10$ ft, $200$ ft$^2$ --- watch for ``$10$'' as the whole answer. Show me
$h=-16t^2+32t+48$, ``how high?'': answer \textbf{2}, $64$ ft. \textbf{1} ($1$ s) = the vertex's
input --- \emph{the target misconception}; \textbf{3} ($3$ s) = the zero; \textbf{4} ($48$ ft) =
the starting height. \textbf{The frame not to rush:} before revealing, ask a \textbf{1} and a
\textbf{2} to explain --- ``is your answer in seconds or in feet?''
\end{multicols}
\end{skillbox}

\boxguard
\begin{skillbox}[Final Round (5 min)]{forestbg}
Four mixed problems, one per round, alone, 5-minute timer: (1) $(2,3)$, $(-1,0)$; (2)
$x^2+1=0$, $D=-4$: none; (3) $h=-16t^2+48t+64$ lands at $t=4$\,s; (4) same ball, how high:
$t=1.5$, $h(1.5)=100$ ft (the crux's twin --- ``$1.5$'' is the input). \textbf{Bonus:} the art
club's hoodies, $R=(30+x)(40-x)=-x^2+10x+1200$: $x=5$, price \$35, revenue \$1225. Students
score themselves out of 4; ask for scores on fingers --- that is your exit read. Say the
\emph{Watch out} aloud as you reveal: the maximum value is the output; when it happens is the
input.
\end{skillbox}

\boxguard
\begin{skillbox}[Homework --- scored, started in class, due the first class after two study halls]{goldbox}
\textbf{Source: DeltaMath online}, set \textbf{2.5 Systems \& Modeling} --- nothing in the
packet. Build the set to cover, in order: linear--quadratic systems by substitution (answers as
ordered pairs, including one with the line not solved for $y$); the number of solutions from a
graph and from the discriminant; projectile models (starting height, maximum height and when,
when it lands); a maximum-area or revenue model. \textbf{Formative check:} the first ``how high
does it go?'' problem --- sort into \emph{output} (correct) / \emph{gave the input (the time)} /
\emph{other feature (the zero or the start)}; if the middle pile is large, open the Unit 2 review
by re-asking the crux. \textbf{Preview:} the Unit 2 test --- the practice test in the packet.
\end{skillbox}

\boxguard
\begin{skillbox}[Active Monitoring --- Watch For]{redbox}
\begin{itemize}[leftmargin=1.2em, nosep]
  \item \textbf{Round 1 ladder and Find the mistake:} $x$-values with no $y$; $y$ paired with
        the wrong $x$.
  \item \textbf{Round 2 Show me:} $D$ taken of the parabola instead of the collapsed equation.
  \item \textbf{Round 3 \rung2, \rung3, Show me:} the vertex's input given as the maximum
        (\emph{the target misconception}).
  \item \textbf{Final Round 4:} ``$1.5$'' for how high.
  \item \textbf{Attitude:} a student parked at \rung1\ with nothing written --- crouch, point at
        the example on the handout, ask for just the first step.
\end{itemize}
Cold-call prompts (after the fingers, never before): ``Is a solution a number or a point?'' \;
``Is your answer in seconds or in feet --- and which did the question ask for?''
\end{skillbox}

% ── Teacher notes — one per component, in the order they are used ──
\begin{teachernote}[Slides]
Pacing: warm-up 5, rounds 13 / 13 / 14, Final Round 5, homework 10 --- 60 minutes. In each round
the definition and example together take \textbf{at most 3 minutes} --- explain further at
desks, not at the board. Reveal a ladder when most of the room has \rung2; the \rung3\ students
keep going. If behind, cut Round~2's Show me (or its \rung3), never Round~3's Show me. Students
write in the handout's notes column, which prints without the answers.
\end{teachernote}

\begin{teachernote}[Homework]
The homework is the DeltaMath set \textbf{2.5 Systems \& Modeling}, online --- not in the packet.
Launch problem~1 aloud in the last 10 minutes (students on their devices), then circulate with
the Watch For list in mind. The first ``how high does it go?'' problem is the formative check.
Due the first class after two study halls; announce the date in class and on TurtleNet.
\end{teachernote}

\end{document}
